\subsubsection{On the Promises of the Neuromorphic Advantage.}
\ac{NC} and \acp{SNN} hold the promise of improving over conventional \acp{ANN} in certain AI applications that require real-time processing of temporally sparse data streams, especially under strict power constraints~\cite{bos_sub-mw_2022, shrestha_efficient_2024}. This concept is broadly referred to as the \textit{Neuromorphic Advantage}~\cite{diamond_comparing_2016,aimone_provable_2021,kudithipudi_neuromorphic_2025}. The expected advantage is rooted in algorithmic properties inherent to \acp{SNN}, such as the sparsity in neuronal activations and connectivity, spatially and temporally sparse spike input encoding or increased expressivity through brain-inspired features like synaptic delays or some forms of online adaptivity.

\subsubsection{The Neuromorphic Simulation-to-Hardware Gap.}
To realize the advantage, SNNs need to be deployed on neuromorphic hardware platforms which exploit the aforementioned algorithmic properties via its hardware architecture and SNN-specific processing elements.
Architectural features that support SNNs include event-driven computation, massive parallelism, co-location of memory and processing (reducing the von Neumann bottleneck), inherent scalability, and novel mixed signal or asynchronous digital circuit designs~\cite{richter_dynap-se2_2024,hoppner_spinnaker_2022,davies_taking_2021}.
Despite the compelling promise of \ac{NC}, the practical and widespread realization of the Neuromorphic Advantage faces hurdles due to the limited availability and accessibility of event-driven neuromorphic accelerators. This scarcity creates a \textit{simulation-to-hardware gap}: while many innovative \ac{SNN} algorithms and \ac{NC} concepts demonstrate considerable promise in simulated environments, their empirical validation, performance characterization, and iterative refinement on actual neuromorphic hardware are severely constrained. This gap, in turn, impedes:
\begin{itemize}
    \item \textbf{HW/SW Co-Design Cycle:} The feedback loop, where insights from hardware performance inform algorithmic refinement and new hardware designs \cite{kudithipudi_neuromorphic_2025}.
    \item \textbf{Algorithmic Innovation:} Without access to real hardware, the exploration of novel SNN architectures, training methods, and optimization techniques is often limited to theoretical constructs or simulations that may not translate well to physical devices.
    \item \textbf{Open-Source Ecosystem Maturation:} A robust open-source software ecosystem lowers the entry barrier to neuromorphic research and development. Without broad access to neuromorphic hardware, the neuromorphic equivalents to programming models, compilers, and interoperable APIs cannot mature effectively~\cite{muir_road_2025}.
    \item \textbf{Hardware-Grounded Benchmarking:} Establishing clear, fair, and widely accepted benchmarks that assess performance on actual neuromorphic hardware (rather than just simulation) becomes exceedingly difficult, making objective comparisons between approaches challenging~\cite{yik_neurobench_2025}.
\end{itemize}
\subsubsection{Related Work.}
Over several decades of neuromorphic hardware acceleration research, the field has recently reached integration levels of billions of neurons and synapses on a single chip. The recent review by Kudithipudi et al.~\cite{kudithipudi_neuromorphic_2025} provides a comprehensive overview of the different architectures of neuromorphic accelerators. However, almost all of these systems remain unavailable to the public, serving to close the simulation-to-hardware gap for only a small number of researchers and developers. While a few systems are commercially available, they are heavily constrained in their programmability to meet strict power and area requirements for commercial viability. To properly close the simulation-to-hardware gap and enable rapid innovation in the public research space, we believe it is necessary to provide an open-source framework that allows the development of neuromorphic hardware accelerators that are both programmable and accessible to a broad community.
Recent related work addresses the acceleration of \ac{DL} programmed \acp{SNN} onto \acp{FPGA}, backed by open-source code repositories including \ac{RTL} descriptions and integration with one \ac{DL} training library out of the neuromorphic open-source ecosystem~\cite{leone_syntzulu_2025,carpegna_spiker_2024,saulquin_modnef_2025}.
At first glance, these works appear to provide a solution to the simulation-to-hardware gap, as they enable the deployment of \ac{SNN} models onto \ac{FPGA} platforms. However, these existing hardware architectures primarily support \ac{SNN} graphs that are constrained by the layered, sequential connectivity patterns inherent to conventional \ac{DL} programming models, making them inadequately prepared for arbitrary, highly recurrent \ac{SNN} graphs that may emerge from future neuromorphic programming paradigms. Furthermore, the aforementioned accelerators implement clock-driven updating mechanisms rather than event-driven processing, limiting their ability to exploit both spatial and temporal sparsity inherent in \ac{SNN} computations.
Overall, there is a lack of open-source, event-driven neuromorphic accelerators that are designed to support arbitrary \ac{SNN} graphs and that can be used to close the simulation-to-hardware gap. This gap hinders the development of neuromorphic hardware accelerators that can fully exploit the Neuromorphic Advantage.

\subsubsection{Our Contribution.}
To address the availability of event-driven, open-source frameworks that provide \ac{SNN} acceleration and thus bridge the simulation-to-hardware gap, we introduce Yet Another Neuromorphic Accelerator (YANA). YANA is an FPGA-based digital SNN accelerator that processes and communicates event-by-event, fully exploiting sparsity in activations, SNN connectivity, and event packet transfer. 
YANA leverages the near-memory computing paradigm and its architecture supports the acceleration of arbitrary SNN graphs, providing flexibility beyond the structures commonly found in the \ac{DL} programming model.

We present hardware resource utilization metrics for a single YANA core and this work's experimental prototype. We also demonstrate how YANA's event-driven architecture can exploit temporal and spatial sparsity to reduce inference latency by deploying a suite of \acp{SNN} with varying pruning levels and input sparsities on the Spiking Heidelberg Digits (SHD) \cite{cramer_heidelberg_2022} dataset.

Furthermore, we release the first public open-source version of the YANA framework for the training, optimization, and deployment of SNNs onto a single YANA core.
\footnote[1]{https://github.com/pachideh/yana-prototype}